\section{The hypothesis family: feasibility and separation}\label{sec:bumps}

We now turn the codewords of \Cref{lem:vg} into a finite family of regression
functions inside the Sobolev ball $\mathcal{W}$ that are pairwise separated in
$\Ltwo$. The amplitude $\omega$ is the single free parameter; it is tuned in
\Cref{thm:main}.

\begin{lemma}[Feasibility and $\Ltwo$ separation]\label{lem:bumps}
Let $g,c_g,C_g$ and $u_\tau$ be as in \Cref{def:bumps}, and let
$\{\tau^{(1)},\dots,\tau^{(M)}\}$ be the codewords of \Cref{lem:vg} with
$M_0=m^d$. Then for every codeword,
\begin{align}\label{eq:sobnorm}
\norm{u_{\tau^{(k)}}}_{\Sob}^2 \;\le\; C_g\,\omega^2,
\end{align}
so that $\omega\le \omega_0 := C_g^{-1/2}$ implies $u_{\tau^{(k)}}\in\mathcal{W}$
for all $k$; and for every pair $k\ne k'$,
\begin{align}\label{eq:sep}
\norm{u_{\tau^{(k)}}-u_{\tau^{(k')}}}_{\Ltwo}^2
\;\ge\; \frac{c_g}{8}\,\frac{\omega^2}{m^{2s}}.
\end{align}
\end{lemma}

\begin{proof}
The proof is two change-of-variables computations: the first bounds the
Sobolev norm of a single bump (giving feasibility after summing disjoint
cells), the second turns the Hamming separation of \Cref{lem:vg} into the
$\Ltwo$ separation Eq.~\eqref{eq:sep}. Throughout, fix the substitution
$y=m(x-x^{(j)})$ on the $j$-th subcube, under which $dx = m^{-d}\,dy$ and $y$
ranges over $(0,1)^d$.

\paragraph{Step 1: Sobolev norm of one bump.}
Fix a cell $j$ and a multi-index $\alpha$ with $|\alpha|=r\le s$. Since
$\partial^{\alpha}\bigl[g(m(x-x^{(j)}))\bigr] = m^{r}\,(\partial^{\alpha}g)(m(x-x^{(j)}))$,
\begin{align*}
\int_{[0,1]^d}\Bigl|\partial^{\alpha}\Bigl(\tfrac{\omega}{m^{s}}\,g(m(x-x^{(j)}))\Bigr)\Bigr|^2 dx
\;=\; & ~ \frac{\omega^2}{m^{2s}}\,m^{2r}\int_{[0,1]^d}\bigl|(\partial^{\alpha}g)(m(x-x^{(j)}))\bigr|^2 dx \\
\;=\; & ~ \frac{\omega^2}{m^{2s}}\,m^{2r}\,m^{-d}\int_{(0,1)^d}\bigl|(\partial^{\alpha}g)(y)\bigr|^2 dy \\
\;=\; & ~ \frac{\omega^2}{m^{2(s-r)}}\,m^{-d}\,\norm{\partial^{\alpha}g}_{\Ltwo}^2 \\
\;\le\; & ~ \omega^2\,m^{-d}\,\norm{\partial^{\alpha}g}_{\Ltwo}^2,
\end{align*}
where the first step differentiates the rescaled bump and pulls out the factor
$m^{r}$, the second step applies the change of variables $y=m(x-x^{(j)})$
together with $\operatorname{supp} g\subseteq(0,1)^d$, the third step collects
powers of $m$, and the last step uses $r\le s$ so that $m^{-2(s-r)}\le 1$.

\paragraph{Step 2: feasibility, Eq.~\eqref{eq:sobnorm}.}
The $M_0$ summands in $u_{\tau^{(k)}}$ have disjoint supports, so all their
mixed partials are orthogonal in $\Ltwo$ and the squared norms add. Hence
\begin{align*}
\norm{u_{\tau^{(k)}}}_{\Sob}^2
\;=\; & ~ \sum_{|\alpha|\le s}\sum_{j=1}^{M_0}\tau^{(k)}_j
        \int_{[0,1]^d}\Bigl|\partial^{\alpha}\Bigl(\tfrac{\omega}{m^{s}}\,g(m(x-x^{(j)}))\Bigr)\Bigr|^2 dx \\
\;\le\; & ~ \sum_{|\alpha|\le s}\,\Nzero{\tau^{(k)}}\,\omega^2\,m^{-d}\,\norm{\partial^{\alpha}g}_{\Ltwo}^2 \\
\;\le\; & ~ \omega^2\,m^{-d}\,M_0\sum_{|\alpha|\le s}\norm{\partial^{\alpha}g}_{\Ltwo}^2 \\
\;=\; & ~ \omega^2\,C_g,
\end{align*}
where the first step expands the definitions Eq.~\eqref{eq:sobolev-norm} and
Eq.~\eqref{eq:bump} and uses disjoint-support orthogonality, the second step
inserts the per-bump bound from Step~1, the third step bounds the Hamming
weight $\Nzero{\tau^{(k)}}\le M_0$, and the last step uses $M_0=m^d$ and the
definition $C_g=\sum_{|\alpha|\le s}\norm{\partial^{\alpha}g}_{\Ltwo}^2$. This
proves Eq.~\eqref{eq:sobnorm}; choosing $\omega\le\omega_0=C_g^{-1/2}$ gives
$\norm{u_{\tau^{(k)}}}_{\Sob}\le 1$, i.e. $u_{\tau^{(k)}}\in\mathcal{W}$.

\paragraph{Step 3: $\Ltwo$ separation, Eq.~\eqref{eq:sep}.}
For one cell $j$, the same change of variables (with $r=0$) gives the $\Ltwo$
mass of a single bump,
\begin{align}\label{eq:cell-mass}
\int_{[0,1]^d}\Bigl|\tfrac{\omega}{m^{s}}\,g(m(x-x^{(j)}))\Bigr|^2 dx
\;=\; \frac{\omega^2}{m^{2s}}\,m^{-d}\,\norm{g}_{\Ltwo}^2
\;=\; \frac{c_g\,\omega^2}{m^{2s+d}}.
\end{align}
The difference $u_{\tau^{(k)}}-u_{\tau^{(k')}}$ is supported on exactly the
cells where $\tau^{(k)}_j\ne\tau^{(k')}_j$, on each of which it equals
$\pm\tfrac{\omega}{m^{s}}g(m(x-x^{(j)}))$. By disjointness of supports,
\begin{align*}
\norm{u_{\tau^{(k)}}-u_{\tau^{(k')}}}_{\Ltwo}^2
\;=\; & ~ \sum_{j=1}^{M_0}\1\!\left\{\tau^{(k)}_j\ne\tau^{(k')}_j\right\}
        \int_{[0,1]^d}\Bigl|\tfrac{\omega}{m^{s}}\,g(m(x-x^{(j)}))\Bigr|^2 dx \\
\;=\; & ~ \rhoH\!\left(\tau^{(k)},\tau^{(k')}\right)\cdot\frac{c_g\,\omega^2}{m^{2s+d}} \\
\;\ge\; & ~ \frac{M_0}{8}\cdot\frac{c_g\,\omega^2}{m^{2s+d}} \\
\;=\; & ~ \frac{c_g}{8}\,\frac{\omega^2}{m^{2s}},
\end{align*}
where the first step uses disjoint-support additivity of the $\Ltwo$ norm, the
second step inserts the per-cell mass Eq.~\eqref{eq:cell-mass} and recognizes the
number of differing coordinates as the Hamming distance, the third step
applies the Varshamov--Gilbert separation $\rhoH\ge M_0/8$ from Eq.~\eqref{eq:vg},
and the last step uses $M_0=m^d$. This is Eq.~\eqref{eq:sep}, completing the proof.
\end{proof}
